\paragraph{Preservation with retained constraints}
When constraints remain at $b=1$, preservation depends on both
safe entry and continued safety.
Let $E_{ij}(s,e')$ and $S_{ij}(s,e',y')$ be the transformed safety
relations with current helper $i$ and next helper $j$.
Assume $E_{00}=\rho_e$, $S_{00}=\rho_s$, and $E_{10}=\mathit{false}$.
All original typed variables and logical monitors remain in $D$;
the next helper is observed before the system responds.
Let $A'$ and $G'$ be generalized B\"uchi objectives whose corresponding
state predicates agree with $A$ and $G$ on $b=0$ states, respectively.
Require $A^R=A'\land\G\F\neg b$ and $G^R=G'$ globally;
the predicates may otherwise be arbitrary on $b=1$ states.
The relations in $01$ and $11$ may otherwise be arbitrary, with
no next-system dependencies in their environmental component.

Let $U$ be the strict safety-winning region for the helper-true
subgame, and let $F$ characterize safe entry to that region:
\begin{align}
 U &= \nu Z.\,\CPre_{11}(Z),\label{eq:helper-U}\\
 F &= \{s\in D\mid\forall e'\in E:\ E_{01}(s,e')\Rightarrow{}\nonumber\\
   &\hspace{25mm}\exists y'\in Y:\ S_{01}(s,e',y')\land(e',y')\in U\}.
 \label{eq:helper-F}
\end{align}
Here $\CPre_{11}$ uses $E_{11}$ and $S_{11}$ with the same typed
quantifiers and dead-end convention as $\CPre$.
Define $W_F$ by replacing $C$ with $F$ in the definition of $W_C$,
including the outer intersection with $F$.

\begin{proposition}[Preservation with retained helper-true constraints]
\label{prop:general-helper}
Under these hypotheses,
\[
 W_R(\cdot,1)=U,\qquad W_R(\cdot,0)=W_F.
\]
Consequently, full original-region preservation holds exactly when
$W\subseteq F$. If the original initializer is preserved and $b$
starts at zero, the transformed initial decision is $I_M(W_F)$.
\end{proposition}
\begin{proof}
At $b=1$, persistence keeps the helper true, so the added assumption
fails and only strict safety remains. The fixed point $U$ supplies safe
responses back into $U$; conversely, states with safety-winning residual
strategies form a post-fixed point of $\CPre_{11}$. Thus $W_R(\cdot,1)=U$.

From a winning helper-false state, each legal flip must enter $U$ safely.
The same holds after every original-legal non-flip prefix, so all reached
states satisfy $F$. Restriction to those histories wins $M_F$.
Conversely, from $s\in W_F$, a winning $M_F$ strategy maintains $F$ until a flip, then
enters $U$ and switches to a uniform safety-winning strategy, available
in this finite, fully observed game. Without a flip, original duties
hold; afterwards safety holds and the added assumption fails. Illegal
inputs release strict duties. This proves $W_R(\cdot,0)=W_F$.

Finally, $W_F\subseteq W\cap F$ gives necessity for preservation.
If $W\subseteq F$, every state reached by an original winning strategy
along original-legal histories lies in $W$, hence in $F$; that strategy
wins $M_F$, proving sufficiency.
The unchanged initializer gives $I_M(W_F)$.
\end{proof}

This is standard safety-game reasoning for a one-way helper,
not a new solving algorithm.
A nontrivial $S_{11}$ can still give $U=D$; even a proper $U$ can
permit every winning active state to enter it. Conversely, the
existence of an immediate $S_{01}$ response does not suffice if
that response cannot reach $U$. For example, use one input value,
a Boolean response $y$, true original initial and justice predicates,
$\rho_e=\mathit{true}$, and $\rho_s=y'$. Then $W=C=D$.
Let $E_{01}=E_{11}=\mathit{true}$,
$S_{01}=\rho_s$, and $S_{11}=\neg y\land\neg y'$.
Every state permits a flip response, but each response reaches
$y=1$, outside $U=\{y=0\}$. Thus $F$ and the active winning region
are empty, despite preserving the original first-flip relation.
In the complete literal case $U=D$ and $F=C$; in the
complete corrected case $U=F=D$.
